\documentclass[12pt]{article}
\usepackage{float}
\usepackage{graphicx}
\usepackage{fontspec}
\usepackage{url}
\usepackage{apacite}
\setmainfont{Times New Roman}
\usepackage[letterpaper,top=2cm,bottom=2cm,left=2cm,right=2cm,marginparwidth=1.5cm]{geometry}

\title{Weekly Assignment 8: Analysing behaviour related to misinformation and/or social engineering and behavioural measures for defending against them}
\author{Truong Khang Tobias Nguyen}
\date{}

\begin{document}

\maketitle

\section{Deception-as-a-Service}
The first concept I choose is Deception-as-a-Service (DaaS) as I think that it might be a niche but very interesting choice of measure for organizational cybersecurity. As most measures we have discussed in class such as security training, strengthening cybersecurity culture, or easing verification or reporting procedures have in common that they seek to build resilience inside an organization, DaaS takes a more proactive approach that puts an organization on the offense.
\newline
Methods such as employing decoys such as fake servers, databases, or user accounts that threat actors can't distinguish from real ones can help to instantly identify them as legitimate users don't have any reason to access them. When an attacker engages in a trap, they can be isolated while their tactics and techniques are recorded, allowing security teams to respond and potentially adapt their security measures for productive environments while turning the tables on attackers and have them waste time on fake infrastructure \cite{stoll1988stalking, acalvio2026deception}.
\newline
One method for organizations to enhance cybersecurity behaviour related to DaaS is to encourage reporting behaviour in employees as organizations can use controlled deception scenarios to evaluate whether employees actually report suspicious activities. For example, by sending employees suspicious messages that actually contain a harmless simulated attack, organizations can measure whether employees recognize the suspicious message, if they report it, how long it will take and if they adhere to given reporting procedures. The theory of planned behaviour states that ones' attitude, subjective norms, and perceived behavioural control influence whether they engage in a certain behaviour - mapped to this example, it could be the employees' intention to report suspicious messages, which is influenced by their attitudes towards reporting, perceived expectations from colleagues, and their confidence in identifying and reporting suspicious mesagges \cite{ajzen2011theory}. By providing feedback on these simulated attacks, organizations can strenghten reporting behaviour and overall cybersecurity behaviour.

\section{Social Engineering}
The second concept I choose is Social Engineering (SE) and I will keep the section shorter than the previous one as we already had quite a few discussions about SE during the course of our studies.
\newline
As established during this course, cybersecurity is not just a question of technical, but also of human nature \cite{indrajit2017social}. On that note, SE techniques such as phishing, vishing, or pretexting are used to gain access to organizational assets without directly attacking their IT systems.
\newline
With this in mind, organizations must reinforce their "human firewall" to deal with non-technical attacks. A straightforward solution is security and awareness training - training employees to recognize SE attacks, how to defend themselves, and how to properly report them can already be an efficient way to safeguard assets and strengthens cybersecurity behaviour by teaching employees from all levels of what threats they can receive directly \cite{aldawood2018educating}.

\bibliographystyle{apacite}
\bibliography{references}

\end{document}